\chapter{Euler Half-Period and the Forced Complex Crossing Pair}
\label{chap:euler-half-period-crossings}

\status{Forced complex crossing pair and its simplicity proved. Additional complex zeros of the scale defect remain open. No SOH-G003 or RH claim is made.}

Chapter 35 established the unique positive crossing of
\[
\Delta_a(u)=X(au)-X(u),\qquad
X(u)=\xi\!\left(\frac{u}{1+u}\right),\qquad a>1.
\]
The reciprocal functional symmetry
\[
X(u)=X(1/u)
\]
contains a second, genuinely complex-geometric statement once the punctured \(u\)-plane is uniformized logarithmically.

Set
\[
u=e^\lambda,\qquad L=\log a,
\]
and define
\[
D_a(\lambda)=\Delta_a(e^\lambda).
\]
Then
\[
\boxed{D_a(-L-\lambda)=-D_a(\lambda)}
\]
and, because \(e^{\lambda+2\pi i}=e^\lambda\),
\[
\boxed{D_a(\lambda+2\pi i)=D_a(\lambda)}.
\]

\section{The logarithmic fixed-point family}

The fixed points of the antisymmetry involution modulo the logarithmic period satisfy
\[
\lambda=-L-\lambda+2\pi i k,
\]
so
\[
\boxed{\lambda_k=-\frac{L}{2}+\pi i k},\qquad k\in\mathbb Z.
\]
At every such point,
\[
D_a(\lambda_k)=-D_a(\lambda_k),
\]
hence
\[
\boxed{D_a(\lambda_k)=0}.
\]
Exponentiating collapses this infinite logarithmic family to two points:
\[
\boxed{u_+=a^{-1/2}},\qquad
\boxed{u_-=-a^{-1/2}}.
\]
Equivalently, these are the two fixed points of the reciprocal scale involution
\[
I_a(u)=\frac{1}{au}.
\]

\section{Euler's identity as the exact half-period bridge}

Successive logarithmic representatives differ by
\[
\lambda_{k+1}-\lambda_k=\pi i.
\]
Euler's identity gives
\[
\boxed{e^{i\pi}=-1},
\]
therefore
\[
\boxed{e^{\lambda_{k+1}}
=e^{\lambda_k}e^{i\pi}
=-e^{\lambda_k}}.
\]
Thus the logarithmic half-period \(\pi i\) exchanges the two forced crossing classes, while the full period \(2\pi i\) returns to the same \(u\)-point. This is an exact phase relation, not an analogy.

For the Uroboros scale \(a=32=2^5\),
\[
L=5\log2,
\qquad
\boxed{\lambda_k=-\frac52\log2+\pi i k},
\]
and
\[
\boxed{u_\pm=\pm\frac1{\sqrt{32}}}.
\]

\section{Simplicity of the forced pair}

Recall the centered entire factorization
\[
\xi\!\left(\frac12+z\right)=F(z^2)
\]
and the positive coefficient theorem
\[
F(w)=\sum_{n\ge0}a_nw^n,\qquad a_n>0.
\]
Therefore
\[
F'(w)>0\qquad(w>0).
\]
In logarithmic coordinates,
\[
X(e^\lambda)
=F\!\left(w(\lambda)\right),
\qquad
w(\lambda)=\frac14\tanh^2\frac\lambda2.
\]
Differentiating \(D_a\), using evenness of \(X(e^\lambda)\) and its \(2\pi i\)-periodicity, yields at every forced crossing
\[
D_a'(\lambda_k)
=2\,\frac{d}{d\lambda}X(e^\lambda)
\Big|_{\lambda=L/2+\pi i k}.
\]
For even \(k\),
\[
w=\frac14\tanh^2\frac{L}{4}>0.
\]
For odd \(k\), the Euler half-period produces
\[
\tanh\!\left(z+\frac{i\pi}{2}\right)=\coth z,
\]
so
\[
w=\frac14\coth^2\frac{L}{4}>0.
\]
Since \(L>0\), the corresponding \(w'(\lambda)\) is nonzero in both parity classes. Hence
\[
F'(w)w'(\lambda)\neq0
\]
and therefore
\[
\boxed{D_a'(\lambda_k)\neq0}.
\]

\begin{theorem}[Forced Euler crossing pair]
For every \(a>1\), the scale defect \(\Delta_a\) has the two symmetry-forced simple zeros
\[
\boxed{u=\pm a^{-1/2}}.
\]
Their logarithmic lifts are
\[
\boxed{\lambda_k=-\frac12\log a+\pi i k},
\]
and adjacent lifts are exchanged by the exact Euler half-turn \(e^{i\pi}=-1\).
\end{theorem}

\section{Open complex frontier}

The theorem proves that the pair \(\pm a^{-1/2}\) is present and simple. It does not prove that
\[
Z(\Delta_a)=\{\pm a^{-1/2}\}.
\]
Additional complex zeros would correspond to nontrivial value collisions between the two arguments of the centered entire function \(F\). Excluding them requires a new global value-separation or injectivity mechanism and cannot be obtained from reciprocal symmetry alone.

\section{Proof firewall}

Proved: the logarithmic crossing family, the two \(u\)-classes, Euler half-period exchange, and simplicity of both forced zeros. Open: additional complex scale-defect zeros, global injectivity of \(F\), PF\(_\infty\), SOH-G003 real-rootedness, and the Riemann Hypothesis.
